

%% ===================================================================
\section{불확실성(Uncertainty) 모델링}
\label{sec:lit-uncertainty}
%% ===================================================================

\chg{관측성이 복구 환경의 가시성을 높여 주더라도, 복구 소요 시간 자체에 내재한 확률적 변동까지 제거하지는 못한다. 이 절에서는 이러한 변동을 설명하는 불확실성 모델링을 일상적 변동을 다루는 가우시안 분포와 극단 사건을 다루는 파레토 분포를 중심으로 검토한다.}

%% -------------------------------------------------------------------
\subsection{가우시안 불확실성 분포}
\label{subsec:gaussian-uncertainty}
%% -------------------------------------------------------------------

불확실성 정량화의 가장 기본적인 도구는 가우시안(정규) 분포이다.
\citet{jaynes2003probability}는 확률 이론의 논리적 기초에 관한 체계적 논의를 통해
중심극한정리(Central Limit Theorem, CLT)의 정보이론적 정당성을 제시하였다.
다수의 독립적인 확률 변수의 합이 가우시안 분포에 수렴한다는 CLT는,
실제 복구 과정에서 발생하는 다양한 소규모 불확실성 요인, \chg{곧 인력 가용성과
장비 상태, 네트워크 지연, 절차적 오류 등의} 복합적 효과가
가우시안 형태로 모델링될 수 있음을 이론적으로 뒷받침한다.

본 연구에서는 일상적 운영 불확실성을 가우시안 분포 \chg{$U^{(G)} \sim N(\mu, \sigma^2)$}로
모델링한다. 이 분포는 대칭적 특성을 가지므로, 복구 시간의 단축과 연장 가능성을
동등하게 반영할 수 있다. 구체적으로 $\DRTO$ 프레임워크에서는
$U_{\text{raw}} \sim N(3, 1)$의 원시 불확실성 값을 생성한 후
$z_U = \kappa \cdot k_U \cdot R_0 \cdot U_{\text{raw}} / (R_{\max} - R_0)$으로
시그모이드 입력값을 산출한다.
\chg{가우시안을 C2 조건에 채택한 근거는 앞서 언급한 중심극한정리에 있다. 복구 과정에서
작용하는 다수의 독립적 소규모 불확실성 요인이 누적될 때 그 합이 가우시안에 수렴하므로,
어느 하나의 지배적 요인 없이 여러 작은 변동이 합산되는 일상적 운영 상황을 가우시안으로
표현하는 것이 이론적으로 타당하다.}
그러나 가우시안 분포는 꼬리 확률(tail probability)이 지수적으로 감소하므로,
대규모 중단이나 극단적 지연과 같은 드문 사건(rare event)의 발생 가능성을
구조적으로 과소평가하는 한계를 지닌다.


%% -------------------------------------------------------------------
\subsection{파레토 분포}
\label{subsec:heavy-tail}
%% -------------------------------------------------------------------

현실 세계의 재해 및 중단 사태는 가우시안 분포가 예측하는 것보다 훨씬 빈번하게
극단적 양상을 보인다. \citet{taleb2007blackswan}은 ``블랙 스완''의 개념을 통해
극도로 낮은 확률이지만 발생 시 막대한 영향을 미치는 사건들이 금융, 기술, 자연재해
영역에서 체계적으로 과소평가되고 있음을 경고하였다.

멱법칙(power-law) 분포는 이러한 극단적 사건을 수학적으로 표현하는 핵심 도구이다.
\citet{newman2005powerlaws}은 자연과학, 사회과학, 기술 시스템에서 발견되는
멱법칙 분포의 보편성을 체계적으로 검토하고, 그 확률밀도함수가
\begin{equation}
  f(x) = \frac{\alpha \, x_m^{\alpha}}{x^{\alpha+1}}, \quad x \geq x_m
  \label{eq:pareto-pdf}
\end{equation}
의 형태를 가지는 파레토(Pareto) 분포로 대표됨을 보였다.
여기서 $\alpha$는 형상 모수(shape parameter) 또는 꼬리 지수(tail index)로서,
분포의 꼬리 두께와 모멘트 존재 조건을 결정하는 핵심 파라미터이다.
$\alpha$가 작을수록 꼬리가 두꺼워져 극단값의 출현 빈도가 높아진다.

%% --- 형상 모수 α의 역할 ---
\subsubsection{형상 모수 $\alpha$의 역할}
\label{subsubsec:pareto-alpha}

파레토 분포의 통계적 성질은 형상 모수 $\alpha$의 값에 따라 질적으로 달라진다.
$n$차 모멘트 $E[X^n]$이 유한하기 위해서는 $\alpha > n$이어야 하며,
\chg{이로부터 다음과 같은 임계 구간이 도출된다. $\alpha$가 1 이하이면 평균조차 발산하여
모든 모멘트가 무한하고, 1과 2 사이이면 평균은 정의되나 분산이 발산한다. 2와 3 사이이면
평균과 분산은 유한하나 왜도가 발산하여 분포 형태가 극도로 비대칭이 되고, 3과 4 사이이면
왜도까지 유한하나 첨도가 발산하여 정규분포 대비 극단값의 빈도가 현저히 높아진다.
$\alpha$가 4를 넘으면 4차 이하의 모든 모멘트가 유한해져 꼬리가 상대적으로 얇아지면서
가우시안에 근사한다.}
\p{[}표~\ref{tab:pareto-alpha-properties}\p{]}는 대표적 $\alpha$ 값에 따른
파레토 분포($x_m = 1$)의 통계적 성질을 정리한 것이다.

\begin{table}[htbp]
  \centering
  \caption{형상 모수 $\alpha$에 따른 파레토 분포의 통계적 성질 ($x_m = 1$)}
  \label{tab:pareto-alpha-properties}
  \small
  \begin{tabular}{ccccc}
    \toprule
    $\alpha$ & 평균 $E[X]$ & 분산 $\text{Var}[X]$ & 왜도 & 첨도 \\
    \midrule
    1 & $\infty$ & $\infty$ & --- & --- \\
    2 & 2.000 & $\infty$ & $\infty$ & --- \\
    \textbf{3} & \textbf{1.500} & \textbf{0.750} & $\boldsymbol{\infty}$ & $\boldsymbol{\infty}$ \\
    5 & 1.250 & 0.104 & 4.649 & 70.800 \\
    10 & 1.111 & 0.015 & 2.811 & 14.828 \\
    \bottomrule
  \end{tabular}
\end{table}

$\alpha$가 증가할수록 평균은 $x_m$에 수렴하고 분산은 급격히 감소하여
분포가 $x_m$ 부근에 집중된다.
반대로 $\alpha$가 감소하면 꼬리가 두꺼워져
극단적으로 큰 값의 출현 확률이 높아진다.
이러한 $\alpha$의 역할은 생존함수(survival function)
$\bar{F}(x) = (x_m/x)^{\alpha}$를 통해 직관적으로 확인된다\chg{. 예컨대 $\alpha = 2$일 때
$\bar{F}(10) = 0.01$인 반면, $\alpha = 5$에서는 $\bar{F}(10) = 10^{-5}$로 극단값의 출현
확률이 1/1{,}000 이하로 감소한다.}

사이버 위험(cyber risk) 분야에서 멱법칙 분포의 적합성은 다수의 실증 연구에 의해
확인되어 왔다. 데이터 유출 규모, 사이버공격에 의한 재무 손실, 시스템 장애 지속 시간 등은
정규분포보다 멱법칙 분포에 의해 더 잘 설명되며,
이는 BCM 맥락에서 극단적 복구 지연 시나리오의 모델링에
파레토 분포가 적합함을 시사한다.
이에 본 연구에서는 멱법칙 분포의 대표적 형태인 파레토 분포를 채택하여
C3 조건의 불확실성을 모형화하며,
이후 본문에서 ``파레토 분포''로 통일 표기한다.

극단값 이론(Extreme Value Theory, EVT)은 이러한 파레토 분포의 극단적 사건을
체계적으로 분석하는 통계적 프레임워크를 제공한다~\citep{coles2001extremes}.
\citet{embrechts1997extremal}은 보험 및 금융 분야에서의 극단값 모델링을 위한
수학적 기초를 수립하였으며, \citet{mcneil2015risk}는 정량적 위험관리에서
EVT의 활용 방안을 구체적으로 논의하였다.
\chg{본 연구는 극단값 이론을 직접적인 추정 기법으로 사용하지는 않으나, 그 핵심 결과는
C3 조건에서 파레토 분포를 채택한 이론적 정당성을 제공한다. 극단값 이론에 따르면 광범위한
분포에서 충분히 높은 임계값을 초과하는 부분이 일반화 파레토 분포로 수렴하므로, 복구 시간의
극단적 꼬리를 파레토 분포로 모형화하는 것은 임의의 선택이 아니라 극단값의 점근적 거동에
부합하는 선택이다. 본 연구가 멱법칙과 극단값 이론에서 가져온 것은 ``극단적 사건의 꼬리는
지수적으로가 아니라 멱법칙으로 감소한다''는 명제이며, 이를 C3 채널의 분포 형태로 직접
반영한 것이다.}

%% --- α = 3 선택 근거 ---
\subsubsection{$\alpha = 3$ 선택 근거}
\label{subsubsec:alpha3-rationale}

본 연구에서는 $\DRTO$ 프레임워크의 조건 \textbf{C3}에서 파레토 분포
\chg{$U^{(P)} \sim 6 \cdot \text{Pareto}(\alpha = 3)$}에 기반한 불확실성 채널을 도입하여,
가우시안 분포(\textbf{C2})가 포착하지 못하는 극단적 복구 지연 시나리오를 모델링한다.
\chg{$\alpha = 3$의 선택은 네 가지 근거에 기반한다. 첫째, $\alpha > 2$이면 평균과 분산이
모두 유한하여 표본 통계량의 안정적 추정과 회귀분석이 가능하다. 둘째, $\alpha \leq 4$이면
첨도가 발산하여 가우시안 분포 대비 극단값이 현저히 빈번하게 출현하는 현실적 시나리오를
표현한다. 셋째, $\alpha = 3$은 유한 분산과 무한 왜도 및 첨도가 공존하는 경계에 위치하여
통계적 분석 가능성과 중꼬리(heavy-tail) 특성 사이의 최적 균형점을 제공한다. 넷째,
사이버 보안 사고의 규모 분포에서 보고되는 $\alpha$ 추정값이 2와 4 사이의 중앙 영역에
해당한다는 실증적 정합성을 갖는다.}


%% --- C2·C3 정량적 비교 및 첨도-꼬리 관계 ---
\subsubsection{가우시안과 파레토 분포의 정량적 비교}
\label{subsubsec:c2c3-comparison}

$\DRTO$ 프레임워크가 채택한 두 불확실성 채널, 즉 가우시안 분포
\chg{$U^{(G)} \sim N(3,\,1)$}(C2)과 파레토 분포 \chg{$U^{(P)} \sim 6 \cdot \text{Pareto}(\alpha=3)$}(C3)은
동일한 기대값($3.0$)을 갖도록 척도를 조정하였으나, 꼬리 영역의 거동에서 질적으로
구별된다. \p{[}표~\ref{tab:c2c3-comparison}\p{]}은 두 분포의 꼬리 성질과 주요 분위수를
비교한 것이다.

\begin{table}[htbp]
\centering
\caption{가우시안(C2)과 파레토(C3) 불확실성 분포의 비교}
\label{tab:c2c3-comparison}
\small
\begin{tabular}{l c c l}
\toprule
\textbf{구분} & \textbf{가우시안 (C2)} & \textbf{파레토 (C3)} & \textbf{의미} \\
\midrule
꼬리 분류 & 경꼬리 (light-tail) & 중꼬리 (heavy-tail) & 감소 속도 기반 성질 분류 \\
꼬리 감소 속도 & $\sim e^{-x^2/2}$ (지수적) & $\sim x^{-\alpha}$ (멱법칙) & 중꼬리가 훨씬 느리게 감소 \\
첨도 & $3$ (기준값) & $\infty$ ($\alpha = 3$) & 꼬리 두께의 수치 지표 \\
분산 & $1.0$ & $27.0$ (약 $27$배) & 변동 폭의 차이 \\
왜도 & $0$ (대칭) & 양수 (우편향) & 극단 지연의 비대칭성 \\
P95 & $4.645$ & $10.29$ (약 $2.2$배) & 상위 $5\%$ 복구 시간 \\
P99 & $5.326$ & $21.85$ (약 $4.1$배) & 상위 $1\%$ 복구 시간 \\
중앙값 & $\approx 3.0$ & $1.56$ & 대부분은 빠르되 소수가 급격히 지연 \\
\bottomrule
\end{tabular}
\end{table}

여기서 첨도와 꼬리 분류는 관련은 있으나 서로 다른 개념임에 유의해야 한다.
첨도는 분포의 꼬리 두께를 하나의 수치로 측정한 지표인 반면, 경꼬리와 중꼬리는
꼬리의 감소 속도에 따른 성질 분류이다. 중꼬리 분포는 첨도가 높은 경향이 있으나,
첨도가 높다고 반드시 중꼬리인 것은 아니다. 예컨대 유한 지지(bounded support)를
갖는 분포도 극단적 이봉 구조를 통해 높은 첨도를 가질 수 있으나 경꼬리에 속한다.
파레토($\alpha=3$)는 첨도가 무한대이므로 중꼬리의 전형적 사례에 해당한다.
파레토 분포의 중앙값($1.56$)이 평균($3.0$)보다 현저히 낮다는 사실은, 대부분의 사건이
빠르게 복구되지만 소수의 극단 사례에서 복구가 급격히 지연되어 평균을 끌어올리는
중꼬리의 구조적 특성을 직접 보여준다. 두 분포의 시각적 비교와 P99 수치는
\figref{fig:pareto_concept}(본문 \secref{sec:result_c3})를 함께 참조하라.

이러한 분포 특성의 차이는 실무 맥락에서도 뚜렷이 구별된다. 가우시안 불확실성은
일상적 운영 변동을 대표하는데, 일반적 라우터 장애의 복구 시간이나 정기적 소프트웨어
업데이트의 롤백, 표준화된 서버 교체 절차처럼 복구 시간이 평균 근처에 집중되고
극단적 편차의 확률이 지수적으로 감소하여 대부분 예측 가능한 범위에서 마무리되는
경우가 여기에 해당한다. 경꼬리 특성상 이러한 사건은 P99에서도 평균 대비
$2.3\sigma$ 이내로 한정된다. 반면 파레토 불확실성은 극단적 사건을 대표하는데,
전사 시스템이 암호화되는 랜섬웨어 공격에서 대부분은 백업 복원으로 수 시간 내에
해결되나 소수의 극단 사례에서는 데이터 복호화와 포렌식 분석으로 복구가 수일까지
연장되며(P99에서 약 $13.2$시간 이상), 지진 후 화재와 정전이 동시에 발생하는 복합
재난이나 핵심 부품 공급사의 공장 화재로 인한 공급망 연쇄 장애 또한 개별 사건
복구 시간의 합을 초과하는 비선형적 지연을 유발하여 꼬리 영역의 복구 부담을 급증시킨다.


%% -------------------------------------------------------------------
\subsection{불확실성 분류 체계의 통합적 이해}
\label{subsec:uncertainty-classification}
%% -------------------------------------------------------------------

본 절에서 다룬 가우시안 분포(C2)와 파레토 분포(C3)는 불확실성을 \textbf{분포 형태}
(즉 \textbf{강도 분포}의 모양) 차원에서 분류한 것이다. 그러나 BCM 운영 환경에서 불확실성 전반을 체계적으로
이해하기 위해서는 \citet{walker2003uncertainty}이 제시한 \textbf{불확실성 단계 분류}
프레임워크와 함께 살펴볼 필요가 있다(\p{[}그림~\ref{fig:uncertainty-levels}\p{]}).
이 프레임워크는 우리가 알고 있는 정도(level of knowledge)에 따라 불확실성을 네 단계로
구분한다. 이는 리스크 평가의 표준 분해(리스크 = 빈도 × 강도)와 대응되며, 분포는 강도
(Severity)에, 시나리오·확률은 빈도(Frequency)에 해당한다. \chg{네 단계는 i) 분포와
확률을 모두 추정할 수 있는 통계적 불확실성, ii) 대안 시나리오는 열거할 수 있으나 확률이
미정인 시나리오 불확실성, iii) 분포 형태 자체를 특정할 수 없는 인지된 무지, iv) 모델과
시나리오와 확률이 모두 불확실한 전적 무지로 구분된다.}

본 연구의 $\DRTO$ 모델은 이 분류 체계 중 양 극단을 정량적으로 다룬다.
통계적 불확실성은 가우시안 분포(C2, \chg{$U^{(G)} \sim N(3,1)$})에 대응하며, 인지된 무지는
파레토 분포(C3, \chg{$U^{(P)} \sim 6 \cdot \text{Pareto}(\alpha\!=\!3)$})로 보수적으로 표현된다.
시나리오 불확실성은 관측성 채널($O_{\text{raw}}$, $k_O$)의 정보 확보를 통해 부분적으로
흡수되며, 전적 무지(블랙 스완 사건)는 본 모델의 적용 한계로 4.3절에서 별도 논의한다.
\chg{한편 운영 단계에서 $U_{\text{raw}}$를 추정하기 위한 데이터 출처(과거 장애 이력 DB,
외부 위협 지수)와 합성 방식은 본문 \subsecref{app:operationalization} 조작화 예시에서 참고용으로 다룬다.}

\begin{figure}[htbp]
\centering
\begin{tikzpicture}[
  every node/.style={text=black},
  level/.style={draw, thick, rounded corners=3pt, text width=125mm,
    minimum height=24mm, font=\large\bfseries, align=left, inner sep=6pt, text=black},
  arr/.style={-{Stealth[length=5mm]}, ultra thick, black}
]
  \node[level, fill=green!15]                   (l1) {Statistical uncertainty (통계적 불확실성)\\
    \normalsize\mdseries 분포(강도)·확률(빈도) 추정 가능. 반복 관측으로 정량화. 예: 장비 고장률.
    \;\textit{$\to$ DRTO C2 가우시안}};
  \node[level, fill=blue!12, below=2mm of l1]   (l2) {Scenario uncertainty (시나리오 불확실성)\\
    \normalsize\mdseries 대안 시나리오 열거 가능, 확률(빈도) 미정. 예: 미경험 재난 유형.
    \;\textit{$\to$ 관측성 채널}};
  \node[level, fill=orange!18, below=2mm of l2] (l3) {Recognized ignorance (인지된 무지)\\
    \normalsize\mdseries 분포(강도) 형태 자체 특정 불가. 예: 신종 사이버공격, 블랙 스완.
    \;\textit{$\to$ DRTO C3 파레토}};
  \node[level, fill=red!12, below=2mm of l3]    (l4) {Total ignorance (전적 무지)\\
    \normalsize\mdseries 모델·시나리오·확률(빈도·강도) 모두 불확실. 예: 기후 복합재난.
    \;\textit{$\to$ 4.3절 한계}};

  \coordinate (atop) at ($(l1.north west) + (-12mm, 0)$);
  \coordinate (abot) at ($(l4.south -| atop)$);
  \draw[arr] (atop) -- (abot);
  \node[font=\normalsize\bfseries, rotate=90, text=black, anchor=south]
    at ($(atop)!0.5!(abot) + (-3mm, 0)$) {불확실성 정도 증가};
\end{tikzpicture}
\caption{불확실성의 단계별 분류 체계와 $\DRTO$ 모델 대응}
\label{fig:uncertainty-levels}
\end{figure}


%% -------------------------------------------------------------------
\subsection{불확실성 하위개념과 업무연속성 리스크와의 관계}
\label{subsec:uncertainty-risk}
%% -------------------------------------------------------------------

\chg{본 연구에서 불확실성은 복구시간의 확률적 변동을 지칭하며, BCM에서
다루는 리스크(재난)와는 분석 수준이 다르다. 리스크는 무엇이 발생하는가라는
사건(event)에 초점을 두는 반면, $\DRTO$에서의 불확실성은 복구시간이 얼마나 변동하는가라는
결과(outcome)에 초점을 둔다. 즉 리스크 사건이 발생하면 복구시간에 불확실성이
유입되고, $\DRTO$가 이를 확률분포로 정량화한다.

경제학자 \citet{knight1921risk}는 리스크와 불확실성을 다음과 같이 구분하였다. 리스크는 결과의
확률분포가 알려져 정량적 측정이 가능한 상황을, 불확실성은 확률분포 자체를 특정할 수
없는 상황을 의미한다. 이 구분은 본 연구의 분류 체계와 대응한다. C2의 가우시안
불확실성은 분포 형태와 매개변수가 사전에 특정되므로 나이트적 의미의 리스크에 가깝고,
C3의 파레토 불확실성은 분포 형태를 확정할 수 없는 나이트적 불확실성을 보수적으로
근사한 것이다. 불확실성을 학술적으로 더 세분한 분류와 본 모델의 대응은
\p{[}표~\ref{tab:app-uncertainty-taxonomy}\p{]}과 같다.}

\begin{table}[htbp]
\centering
\caption{불확실성 하위개념 분류와 $\DRTO$ 모델 대응}
\label{tab:app-uncertainty-taxonomy}
\small
\begin{tabular}{p{1.8cm} p{2.0cm} p{3.8cm} p{2.8cm} p{2.8cm}}
\toprule
\textbf{분류} & \textbf{영어 표기} & \textbf{정의} & \textbf{BCM 예시} & \textbf{DRTO 대응} \\
\midrule
인식론적 불확실성 & Epistemic & 지식 부족에 기인, 정보 확보로 축소 가능 & 미경험 재난 유형, 불완전한 BIA & 관측성 채널 ($O_{\text{raw}}$, $k_O$) \\
\midrule
우연적 불확실성 & Aleatory & 본질적 무작위성, 데이터를 추가해도 축소 불가 & 장비 고장률, 네트워크 지연 변동 & C2 가우시안 ($U^{(G)}$) \\
\midrule
나이트적 불확실성 & Knightian & 확률분포 자체를 특정할 수 없는 불확실성 & 블랙 스완 사건, 신종 사이버공격 & C3 파레토 ($U^{(P)}$) \\
\midrule
심층 불확실성 & Deep Uncertainty & 모델·시나리오·확률 모두 불확실 & 기후변화 복합재난, 팬데믹 & 한계 사항 (\secref{sec:limitations}) \\
\bottomrule
\end{tabular}
\end{table}

\chg{표에서 우연적 불확실성은 반복 관측으로 분포를 추정할 수 있으나 변동 자체를 제거할
수는 없는 유형이고, 인식론적 불확실성은 관측성($O_{\text{raw}}$)을 높임으로써 축소할 수
있어 $\DRTO$ 모델의 관측성 채널이 직접 반영하는 메커니즘에 해당한다. 나이트적 불확실성은
분포 형태를 사전에 특정할 수 없으므로 본 연구에서는 파레토 분포로 보수적으로 모델링하며,
심층 불확실성은 현 모델의 범위를 벗어나 \secref{sec:limitations}에서 한계로 명시하였다.}

\chg{한편 BCM에서 다루는 세 가지 재난 유형과 본 모델의 불확실성은
\p{[}표~\ref{tab:app-disaster-mapping}\p{]}과 같이 대응한다. 동일한 재난 유형이라도
발생 규모에 따라 불확실성 특성이 달라져, 일상적 서버 장애는 가우시안(C2)으로,
대규모 사이버공격은 파레토(C3)로 모형화된다.}

\begin{table}[htbp]
\centering
\caption{BCMS 재난 유형과 $\DRTO$ 불확실성 매핑}
\label{tab:app-disaster-mapping}
\small
\begin{tabular}{p{2cm} p{3.2cm} p{2.8cm} p{2.5cm} p{2.5cm}}
\toprule
\textbf{재난 유형} & \textbf{세부 예시} & \textbf{빈도-영향 특성} & \textbf{불확실성 유형} & \textbf{DRTO 조건} \\
\midrule
자연재난 & 지진, 태풍, 홍수, 산사태 & 빈도 낮음, 영향 극대 & 나이트적 & C3 (파레토) \\
\midrule
사회재난 & 대규모 정전, 통신장애, 공급망 단절, 테러 & 빈도 중간, 영향 중$\sim$대 & 우연적 + 나이트적 & C2 $\sim$C3 \\
\midrule
기술적 재난 & 서버 장애, 사이버공격, 소프트웨어 결함, 데이터 유출 & 빈도 높음, 영향 소$\sim$극대 & 우연적(일상) / 나이트적(극단) & C2(일상) / C3(극단) \\
\bottomrule
\end{tabular}
\end{table}

\chg{사회재난의 경우에도 공급망 단절의 파급 범위에 따라 가우시안(C2)에서 파레토(C3)로
전이될 수 있다.}

\chg{끝으로 BCM에서의 리스크와 $\DRTO$에서의 불확실성의 유사점과 차이점을
\p{[}표~\ref{tab:app-risk-vs-uncertainty}\p{]}에 비교하였다. 리스크가 사건의 발생
가능성과 영향을 다루는 데 비해 불확실성은 그 결과인 복구시간의 확률적 변동을 다루며,
양자는 모두 복구시간에 영향을 미치므로 $\DRTO$가 이를 통합 정량화한다.}

\begin{table}[htbp]
\centering
\caption{BCMS 리스크(재난)와 $\DRTO$ 불확실성 비교}
\label{tab:app-risk-vs-uncertainty}
\small
\begin{tabular}{p{2.5cm} p{5.5cm} p{5.5cm}}
\toprule
\textbf{비교 기준} & \textbf{BCMS 리스크 (재난)} & \textbf{DRTO 불확실성} \\
\midrule
정의 & 조직 목표에 부정적 영향을 미치는 사건의 발생 가능성과 결과~\nrev{(ISO 31000, 2018)} & 복구시간($\DRTO$)의 확률적 변동 범위 \\
\midrule
분석 단위 & 사건(event): ``무엇이 발생하는가'' & 결과(outcome): ``복구시간이 얼마나 변동하는가'' \\
\midrule
방향성 & 주로 부정적 (위협 중심) & 양방향 (단축도 연장도 가능) \\
\midrule
정량화 방법 & 발생가능성 $\times$ 영향도 매트릭스 & 확률분포 (가우시안 / 파레토) \\
\midrule
시간적 성격 & BIA 시점 고정 평가 (정적) & 실시간 동적 변동 \\
\midrule
축소 방법 & 예방·완화 통제 수단 & 관측성 향상 ($O_{\text{raw}}$ $\uparrow$) \\
\midrule
\multicolumn{3}{l}{공통점: 둘 다 복구시간에 영향 $\to$ $\DRTO$가 양자를 통합 정량화} \\
\bottomrule
\end{tabular}
\end{table}

\chg{이상에서 살펴본 가우시안과 파레토 불확실성은 모두 복구 시간에 확률적 변동을 유발하므로, 이를 $\DRTO$ 모델에 반영하려면 그 영향을 일정한 경계 안에서 변환하는 장치가 필요하다. 다음 절에서는 이러한 변환을 담당하는 시그모이드 함수와 바운딩 기법을 검토한다.}
